\chapter{Conclusion and Future Work}
\label{chap:conclusion}

\section{Introduction}

This chapter concludes the NAJDA graduation project by summarizing the work presented throughout this thesis. It reviews the project objectives, highlights the major achievements, discusses the limitations of the current implementation, and outlines potential directions for future enhancement.

The conclusions presented in this chapter are based on the design, implementation, and evaluation activities described throughout the previous chapters and follow accepted software engineering practices for project assessment and continuous improvement \cite{pressman2020se,sommerville2016software}.

%------------------------------------------------

\section{Project Summary}

NAJDA was developed as an emergency dispatch simulation platform intended for academic and educational purposes. The project demonstrates how modern software engineering practices can be applied to simulate emergency response operations without requiring integration with governmental emergency infrastructures.

The platform supports multiple user roles, including citizens, dispatchers, ambulance crews, firefighters, police officers, hospital staff, and system administrators. By combining real-time communication, distributed software architecture, and artificial intelligence-assisted decision support, NAJDA provides a realistic representation of emergency dispatch workflows while maintaining human control over operational decisions. The adopted architecture demonstrates how modern software engineering principles can be integrated to produce scalable and maintainable distributed applications \cite{bass2022software,newman2021microservices}.

%------------------------------------------------

\section{Objectives Achievement}

The primary objectives established at the beginning of the project have been addressed through the design and implementation of the proposed platform.

The project successfully demonstrates:

\begin{itemize}
    \item A realistic emergency dispatch simulation workflow.
    \item Multi-role coordination between emergency participants.
    \item Event-driven communication between distributed software components.
    \item AI-assisted decision support while preserving human authority.
    \item Real-time incident monitoring and mission management.
    \item A scalable and modular software architecture suitable for future expansion.
\end{itemize}

\textit{Placeholder: Final evaluation of objectives after project completion.}

%------------------------------------------------

\section{Project Contributions}

The NAJDA platform contributes to software engineering education by providing a comprehensive case study that combines multiple modern technologies within a single integrated application.

The main contributions include:

\begin{itemize}
    \item Design of an event-driven microservice architecture for emergency dispatch simulation.
    \item Integration of multiple emergency service roles within a unified platform.
    \item Application of artificial intelligence as a decision-support mechanism rather than an autonomous decision-maker.
    \item Demonstration of real-time communication techniques for mission coordination.
    \item Development of an extensible platform that can support future academic research and educational activities.
\end{itemize}

%------------------------------------------------

\section{Project Limitations}

Although NAJDA demonstrates the core concepts of a modern emergency dispatch platform, several limitations remain.

\begin{itemize}
    \item The platform operates solely as an academic simulation.
    \item Integration with governmental emergency communication infrastructure is intentionally excluded.
    \item Artificial intelligence models rely on limited datasets suitable for educational purposes.
    \item Performance evaluation depends on the available development and testing environment.
    \item Certain advanced routing and traffic prediction capabilities remain outside the scope of the current implementation.
\end{itemize}

%------------------------------------------------

\section{Future Work}

Several improvements may be incorporated in future versions of NAJDA to enhance both realism and functionality.

Potential future work includes:

\begin{itemize}
    \item Integration with real-time traffic and routing services.
    \item Advanced geographic information system (GIS) capabilities.
    \item Improved emergency demand prediction using larger datasets.
    \item Enhanced AI models trained on real emergency response data.
    \item Offline communication mechanisms for unstable network environments.
    \item Multi-region deployment with distributed infrastructure.
    \item Expanded analytics and operational reporting dashboards.
    \item Wearable device integration for emergency responders.
    \item Tablet-optimized dispatcher interfaces.
    \item Comprehensive usability studies involving emergency management professionals.
\end{itemize}

%------------------------------------------------

\section{Lessons Learned}

The development of NAJDA provided valuable practical experience in software engineering, distributed systems, artificial intelligence, cloud computing, mobile application development, and system integration. The project also demonstrated the importance of requirements engineering, modular architecture, effective communication between software components, and iterative development throughout the software development lifecycle.

\textit{Placeholder: Add team-specific lessons learned after project completion.}

%------------------------------------------------

\section{Final Remarks}

NAJDA demonstrates that a realistic emergency dispatch platform can be developed as an educational simulation while incorporating modern software engineering principles, distributed system architecture, and artificial intelligence-assisted decision support. The project serves as both a practical software engineering solution and an academic case study illustrating the design and implementation of complex real-time information systems.

It is anticipated that the concepts, architecture, and implementation presented in this thesis will provide a solid foundation for future research, educational projects, and further development in intelligent emergency management systems.

%------------------------------------------------

\section{Chapter Summary}

This chapter concluded the NAJDA project by summarizing the work presented throughout the thesis, evaluating the achievement of the project objectives, discussing current limitations, and identifying opportunities for future enhancement. Together with the previous chapters, it completes the documentation of the proposed emergency dispatch simulation platform.